\documentclass[10pt,twocolumn]{article}
\usepackage[scaled]{helvet}
\renewcommand\familydefault{\sfdefault}
\usepackage[T1]{fontenc}
\usepackage[margin=0.7in]{geometry}
\usepackage{titlesec}
\usepackage{enumitem}
\usepackage{parskip}
\setlength{\columnsep}{0.24in}
\setlist[itemize]{leftmargin=1.2em, topsep=0.2em, itemsep=0.18em}
\titleformat{\section}{\large\bfseries}{\thesection}{1em}{}

\begin{document}
\title{\vspace{-1.6cm}AWS API Gateway}
\author{AWS ML Study Notes}
\date{}
\maketitle
\small

\section*{Overview}
API Gateway is a managed front door for HTTP, REST, and WebSocket APIs. It receives client traffic, applies policies, and forwards requests to backend integrations.

For ML applications, API Gateway often sits in front of inference services, orchestration Lambdas, or thin application layers that expose model capabilities to external consumers.

\section*{Core Concepts}
\begin{itemize}
\item Core features include routing, request validation, throttling, authorization, custom domains, stage management, and integration mapping.
\item Backends can be Lambda functions, HTTP services, private integrations, or other AWS services depending on the API design.
\item Usage plans, API keys, authorizers, and logging help separate public access patterns from internal or partner access patterns.
\end{itemize}

\section*{How It Works}
\begin{itemize}
\item A client sends a request to the API endpoint, API Gateway applies authentication and request handling rules, then forwards the request to the configured integration.
\item The backend returns a response, which API Gateway can transform, enrich with headers, or standardize before returning to the caller.
\item Operational features such as throttling and stage based rollout help protect backends and manage gradual release of new behavior.
\end{itemize}

\section*{AI and ML Relevance}
\begin{itemize}
\item ML APIs benefit from API Gateway because it centralizes security, quotas, and routing policies before traffic reaches expensive inference components.
\item It is also a clean way to expose asynchronous prediction patterns or file based workflows where the initial request starts work and later calls fetch results.
\end{itemize}

\section*{Operational Notes}
\begin{itemize}
\item Be aware of payload sizes, timeout limits, and latency overhead. Not every inference pattern fits cleanly through a synchronous API boundary.
\item Use structured request validation and authorization rather than trusting each backend to implement every guard consistently.
\item Logging and correlation identifiers matter. When an inference path spans API Gateway, Lambda, and a model backend, debugging needs end to end traceability.
\end{itemize}

\section*{What To Remember}
\begin{itemize}
\item API Gateway is the contract boundary, not the business logic itself.
\item It protects and shapes backend access; it does not replace backend capacity planning.
\item In ML systems, think about cost per request and latency amplification before exposing a heavy model behind a public API.
\end{itemize}

\end{document}
